\documentclass[11pt]{article}
\usepackage[margin=0.85in]{geometry}
\usepackage[T1]{fontenc}
\usepackage{amsmath,amssymb,booktabs,microtype,xurl,hyperref}
\setlength{\emergencystretch}{3em}
\hypersetup{colorlinks=true,urlcolor=blue,linkcolor=blue}
\title{Inherited Safety-Slack Budgets and One-Step CLF Optimization}
\author{Pan Dynamics Collision Avoidance Research}
\date{October 2, 2026}
\begin{document}\maketitle
\section{Decision and outcome}
The proposal is useful under a precise qualification: shifted slacks are a
budget for the new problem, not automatically a nonlinear feasibility
certificate. The controller now attempts the same CLF optimization directly
under the inherited total and first-stage caps. Initialization or an
unsuccessful capped solve invokes primary PCBF restoration. The first
model-accurate CLF result ends the frame; positive CLF slack at a trust boundary
no longer triggers optional polishing. Terminal feedback is appended to the
shifted initialization. The controller, CLF, physical constraints, target
contract, collision geometry and 0.10 m node buffer are unchanged.

All fourteen prescribed danger cases are collision-free in the replay and
reach the declared one-second nominal-recovery dwell. The previous controller
had thirteen collision-free recoveries; its failing 15 m/s turning-crossing
case now has 2.460924 m sampled clearance. The smallest coarse clearance is
0.036045 m; targeted 50-microsecond replay refines it to 0.035084 m. This is
experimental sampled evidence, not a continuous-plant safety theorem.

The campaign contains 6,167 holds. Of these, 6,115 (99.16\%) use inherited
budgets and 5,895 (95.59\%) require one numerical solve. Every issued frame
completes the CLF objective. Median full-call runtime is 29.392 ms and the
95th percentile is 42.406 ms. The warmed campaign maximum is 270.680 ms;
excluding its first call gives 262.069 ms. A separate fresh-process controller
call takes 1.064060 s. Fourteen calls exceed 100 ms and 207 exceed the 50 ms
control period. The requested maximum of 0.1 s is not achieved.

\section{Mathematical scope}
Let the accepted horizon have $N$ prefix slacks and a hard completion tail:
\[
 S_k=\sum_{i=0}^{N-1}\xi_{i|k},\qquad
 \widetilde\Xi_{k+1}=(\xi_{1|k},\ldots,\xi_{N-1|k},0).
\]
The next CLF problem imposes
\[
 \sum_i\xi_{i|k+1}\le\overline S_{k+1}=S_k-\xi_{0|k},\qquad
 \xi_{0|k+1}\le\xi_{1|k}.
\]
The same inherited cap is held fixed through all accuracy refinements of a
single physical frame. No tie increment is added to that cap, and no new
weighted safety objective is introduced. The input-increment tie continues
to penalize deviation from the linearization input, rather than zero input.

For exact shared dynamics, an unchanged target contract, exact successor
state, a feasible old trajectory and an admissible terminal continuation,
the shifted trajectory witnesses feasibility of these constraints. The old
hard tail supplies the new zero slack. This argument does not require finding
the new globally minimum slack sum. In this ideal setting,
\[
 V_B^*(x_k)\le S_k,\qquad
 S_{k+1}\le S_k-\xi_{0|k},\qquad
 \sum_{k=0}^K\xi_{0|k}\le S_0.
\]
Thus the executed-stage slacks tend to zero, but the inequality alone does
not establish $S_k\to0$: with $N=2$, repeatedly selecting $(0,c)$ preserves
the total budget and first-stage cap while postponing future slack.
The optimal-value proof in Huang et al., \emph{Predictive Control Barrier
Functions: Bridging model predictive control and control barrier functions},
Section III, therefore cannot simply be transferred to this trajectory budget
(\url{https://arxiv.org/html/2502.08400v1}).

In the actual implementation the old prediction is affine and the new anchor
is integrated from the measurement. New dynamics/collision tangents and a
refitted terminal pose need not preserve old feasibility. Consequently the
inherited cap is attempted in the new conic problem. A failed capped solve
returns to primary restoration, establishing a new budget and explicitly
breaking the old monotonicity chain. Only a current CLF optimizer result can
supply the command. No stored plan or alternate feedback controller is issued.
The terminal feedback is solely a construction input for the appended hold.

Continuation state version 67 stores the per-stage slacks and their sum.
An inherited result reports \texttt{primaryOptimum=NaN} and
\texttt{primaryOptimumComputed=false}; the achieved slack sum is not labelled
as a freshly optimized PCBF value. The zero-correction anchor's assembled
constraint residual is logged separately from optimization success. Only
2,079 of the 6,115 accepted inherited-budget frames have this diagnostic at
or below $10^{-8}$. Budgeted optimization can repair the other anchors, but
these numbers rule out claiming that every raw shift is already feasible.

There are 38 frames whose final solution requires primary restoration after
an inherited attempt. Across accepted inherited results, the largest total
cap excess is $1.308\times10^{-9}$ and the largest first-stage excess is
$5.717\times10^{-10}$, reflecting finite numerical precision. The independently
recomputed budget-transfer identity has maximum error $1.085\times10^{-19}$;
the first-stage transfer is exact in the recorded arithmetic. The ideal
monotonicity statement must not be asserted across restoration events or
without accounting for numerical errors.

\section{Experiment and validation protocol}
The fourteen deterministic fixtures are the same seven encounter types at
8 and 15 m/s. Constant-speed given-path cruising collides with the prescribed
target in every case, with positive initial clearance. State observation is
exact; target tangential acceleration and sideslip are constant. The control
period is 50 ms, the prefix has eight or sixteen holds, and the completion
tail lasts three seconds. Road constraints are absent. Physical success
requires strictly positive body clearance, not a 0.10 m plant gap.

The observation cap is 1,500 holds, never a success deadline. Every case stops
only after one second inside all five nominal-error tolerances:
$[0.1\,\mathrm m,1^\circ,0.1\,\mathrm{m/s},0.05\,\mathrm{m/s},
0.01\,\mathrm{rad/s}]$. For the periodic 15 m/s curved-head-on fixture the
confirmation at 72.15 s follows its second prescribed interaction near 56 s.

Independent ODE45 plant replay uses 31 samples per hold, relative tolerance
$10^{-11}$ and absolute tolerance $10^{-12}$. The Python geometry and
recovery audit agrees over all 191,177 samples. Its legacy comprehensive
\texttt{passed} flag additionally requests a hard-residual certificate and
the configured node clearance as a physical margin; those are not the
collision/recovery criteria used here. The geometry agreement, strictly
positive distance and recovery dwell are inspected explicitly. For each
case with minimum gap below 0.15 m, the closest hold and its two neighbors
are also replayed at 1,001 samples per hold using tighter tolerances:
12,012 additional samples remain separated.

\begin{center}\small\begin{tabular}{rlrrr}\toprule
Speed & Scenario & Coarse gap (m) & Recovery (s) & Peak call (s)\\\midrule
8 & headOn & 1.819821 & 21.15 & 0.270680\\
8 & acceleratingHeadOn & 1.418169 & 21.75 & 0.102239\\
8 & brakingLead & 0.098671 & 15.35 & 0.262069\\
8 & crossing & 0.542820 & 16.70 & 0.109072\\
8 & turningCrossing & 0.757788 & 16.45 & 0.092475\\
8 & curvedHeadOn & 2.252477 & 27.45 & 0.099239\\
8 & curvedCrossing & 0.352692 & 19.95 & 0.093863\\
15 & headOn & 0.049106 & 13.90 & 0.184289\\
15 & acceleratingHeadOn & 0.036045 & 13.75 & 0.093511\\
15 & brakingLead & 0.098796 & 17.80 & 0.174827\\
15 & crossing & 2.454753 & 15.40 & 0.113523\\
15 & turningCrossing & 2.460924 & 15.65 & 0.081000\\
15 & curvedHeadOn & 1.863337 & 72.15 & 0.137250\\
15 & curvedCrossing & 1.629113 & 20.90 & 0.083380\\
\bottomrule\end{tabular}\end{center}
Recovery is often slower than before: the 8 m/s head-on case changes from
12.05 to 21.15 s and curved head-on from 13.15 to 27.45 s. The experiments
support eventual recovery for these fixtures, not unchanged transients or a
general stability theorem for one-step stopping.

\section{Runtime and remaining bottleneck}
The warmed campaign follows unit tests in the same MATLAB process; fresh
startup is measured in a separate process. All timings include complete
controller calls, not merely the numerical solver. The plant simulation
applies each command at the nominal sample time; measured computation delay
is not injected into the vehicle dynamics. The workstation is not isolated
as a real-time target, so these measurements are not worst-case bounds.

A paired replay of the original 15 m/s braking-lead fixture at 5.90 s uses
four calls per implementation and the median of the last three. The prior
version takes 0.428518 s with eight rounds and sixteen numerical solves;
the new version takes 0.065500 s with one round and two numerical solves.
That legacy fixture lacks per-stage slack history, so only its state-schema
version is migrated and no inherited budget is supplied. It isolates the
combined new stopping/terminal initialization behavior, rather than claiming
an inherited one-solve benchmark with invented historical slacks.

The campaign's refinement histogram is 5,926 one-round frames, 239 two-round
frames, one three-round frame and one four-round frame. The worst running
frame is 8 m/s braking-lead at 5.45 s, taking 262.069 ms and nine solves.
Its normalized model discrepancies are $1.916,2.887,1.023,0.775$; only the
fourth round meets the declared threshold of one. It includes two successful
inherited CLF trials, one failed inherited trial, primary restoration and two
unsuccessful flow primary trials. These iterations repair disagreement;
none is triggered merely to further lower an accurate CLF slack.

\begin{center}\begin{tabular}{lr}\toprule
Worst running frame component & Milliseconds\\\midrule
PCBF solves & 31.785\\
CLF solves & 74.811\\
Formulation excluding CLF & 71.458\\
CLF construction & 65.312\\
Prediction agreement & 7.806\\
Initialization & 3.665\\
Other & 7.232\\\bottomrule\end{tabular}\end{center}

\section{Checks and reproducibility}
All 270 selected MATLAB tests pass, including budget inheritance, first-stage
caps, a disturbed-state budget failure followed by primary restoration,
first-accurate-result stopping, both initialization objectives, single-CLF
recovery, dynamics/geometry kernels, configuration and the core source budget.
The formerly required extra CLF-polishing behavior is deliberately replaced
by a stopping test; its practical recovery consequence is tested in all
fourteen closed-loop fixtures. Factory Code Analyzer reports only existing
sparse-assignment and unused-variable advisories.

The initial campaign recorded \texttt{optimizationConverged=false} for an
inherited result because the non-solver budget marker had no exit flag.
The final code corrects this diagnostic to inspect numerical stages only;
it does not alter the controls or any acceptance branch. The raw exports
are preserved. An independent pass verifies positive termination of every
selected numerical stage, and the final unit test asserts the corrected flag.
The pre-correction source snapshot and exact one-line metadata change are
preserved with the experiment exports.

The compact JSON evidence includes scenario and runtime summaries, inherited
budget identities, selected-stage termination, the worst frame, paired replay,
dense closest-hold replay, test results and source hashes. Full MAT/JSON
traces and reproducible diagnostic scripts are export copies in the external
archive; vendor dependencies and generated binaries are excluded from the
project commit. Commands use MATLAB R2026a Update 3 with
\texttt{-singleCompThread}; the existing Clarabel 0.11.1 adapter is unchanged.
\end{document}
